\section{Use case and discussion}
PolyArt is designed to support a broad range of downstream research tasks in multimodal art understanding.The four independently generated perspectives enable perspective-conditioned fine-tuning: researchers can train or evaluate vision-language models on each dimension separately --- formal, emotional, historical, or narrative --- and measure the effect of perspective alignment on generation quality, without the confounds introduced by mixing heterogeneous annotation styles.
The unified caption ($e^{\text{unified}}$) provides a single rich descriptor for cross-modal retrieval experiments, where multi-perspective grounding produces more semantically complete query representations than any single-perspective caption.
The co-presence of ArtEmis emotion labels and generated emotional perspectives makes PolyArt directly applicable to affective computing research on visual art, supporting both classification and generation tasks with human-grounded supervision.
Finally, the ArtRAG-grounded historical perspectives, combined with style and artist metadata, enable knowledge-intensive art question answering and retrieval-augmented generation benchmarks that require models to integrate visual evidence with external cultural knowledge.
Datasets, generation scripts, and pre-computed embeddings are released publicly to facilitate reproducible research across all these applications.

